% ===== PRÉAMBULE CANONIQUE — à coller VERBATIM en tête de CHAQUE .tex =====
\documentclass[11pt,a4paper]{article}
\usepackage[utf8]{inputenc}\usepackage[T1]{fontenc}\usepackage{lmodern}
\usepackage[french]{babel}
\usepackage{amsmath,amssymb,mathtools}\usepackage{icomma}
\usepackage[version=4]{mhchem}          % \ce{...} : équations & formules
\usepackage{siunitx}\sisetup{locale=FR} % \qty{}{}, \si{}, \ang{}
\usepackage{chemfig}                    % \chemfig{...} : structures développées
\usepackage{xcolor}\usepackage{colortbl}
\usepackage{tikz}\usepackage{pgfplots}\pgfplotsset{compat=1.18}
\usepackage[most]{tcolorbox}\usepackage{enumitem}\usepackage{array,booktabs}
\usepackage[a4paper,margin=2.2cm]{geometry}
\usetikzlibrary{arrows.meta,calc,angles,quotes,positioning,patterns,backgrounds,shapes.geometric}
\definecolor{primary}{RGB}{222,49,99}\definecolor{primarydark}{RGB}{150,22,55}
\definecolor{defcol}{RGB}{0,114,178}\definecolor{methcol}{RGB}{52,92,148}
\definecolor{excol}{RGB}{0,135,90}\definecolor{attncol}{RGB}{200,40,55}
\tcbset{boxbase/.style={enhanced,breakable,boxrule=0.9pt,arc=2.5pt,
  left=3mm,right=3mm,top=2.3mm,bottom=2.3mm,fonttitle=\bfseries,coltitle=white}}
% --- boîtes TUTORIEL (noms EXACTS, ne pas renommer) ---
\newtcolorbox{cle}[1][]{boxbase,colback=defcol!6,colframe=defcol,
  title={\textsc{À retenir}\ifx\relax#1\relax\relax\else\textmd{ : \itshape #1}\fi}}
\newtcolorbox{methode}[1][]{boxbase,colback=methcol!7,colframe=methcol,sharp corners,
  title={\textsc{Méthode}\ifx\relax#1\relax\relax\else\textmd{ --- \itshape #1}\fi}}
\newtcolorbox{exemplebox}[1][]{boxbase,colback=excol!5,colframe=excol,
  title={\textsc{Exemple}\ifx\relax#1\relax\relax\else\textmd{ : \itshape #1}\fi}}
\newtcolorbox{piege}[1][]{boxbase,colback=attncol!6,colframe=attncol,
  title={\textsc{Piège}\ifx\relax#1\relax\relax\else\textmd{ : \itshape #1}\fi}}
\newtcolorbox{remarque}[1][]{boxbase,colback=black!4,colframe=black!45,
  title={\textsc{Remarque}\ifx\relax#1\relax\relax\else\textmd{ : \itshape #1}\fi}}
% --- boîtes EXERCICE / CORRIGÉ (noms EXACTS, auto-comptées) ---
\newtcolorbox[auto counter]{exo}[1][]{boxbase,colback=primary!4,colframe=primary,
  title={\textsc{Exercice}~\thetcbcounter\ifx\relax#1\relax\relax\else\textmd{ --- \itshape #1}\fi}}
\newtcolorbox[auto counter]{corrige}[1][]{boxbase,colback=excol!4,colframe=excol,
  title={\textsc{Corrigé}~\thetcbcounter\ifx\relax#1\relax\relax\else\textmd{ --- \itshape #1}\fi}}
\newcommand{\R}{\mathbb{R}}\newcommand{\N}{\mathbb{N}}\newcommand{\Z}{\mathbb{Z}}\newcommand{\ds}{\displaystyle}
% (un \usepackage{fancyhdr}+en-tête est possible ; il est ignoré côté web)
% =========================================================================
\begin{document}
\sisetup{per-mode=symbol}
\begin{center}\small\itshape
Aux TPN : $V_{\mathrm{m}}=\qty{22.4}{\litre\per\mole}$.
\end{center}

\begin{exo}[du sulfure de fer : partir de masses]
Pour $\ce{Fe + S -> FeS}$, on chauffe \qty{5.6}{\gram} de fer avec \qty{6.4}{\gram} de soufre.
On donne $M(\ce{Fe})=\qty{56}{\gram\per\mole}$, $M(\ce{S})=\qty{32}{\gram\per\mole}$,
$M(\ce{FeS})=\qty{88}{\gram\per\mole}$.
\begin{enumerate}[label=\alph*)]
  \item Calcule les quantités de matière de \ce{Fe} et de \ce{S}.
  \item Quel est le réactif limitant ?
  \item Quelle masse de \ce{FeS} se forme ?
\end{enumerate}
\end{exo}

\begin{exo}[limitant en phase gazeuse]
Pour $\ce{2 CO + O2 -> 2 CO2}$ (mêmes température et pression pour tous les gaz), on mélange
\qty{4}{\litre} de \ce{CO} et \qty{3}{\litre} de \ce{O2}.
\begin{enumerate}[label=\alph*)]
  \item Quel est le réactif limitant ? (les volumes se comparent comme des moles)
  \item Quel volume de \ce{CO2} obtient-on ?
  \item Quel volume de \ce{O2} reste-t-il ?
\end{enumerate}
\end{exo}

\begin{exo}[attaque du zinc par l'acide]
Pour $\ce{Zn + 2 HCl -> ZnCl2 + H2 ^}$, on plonge \qty{6.5}{\gram} de zinc
($M=\qty{65}{\gram\per\mole}$) dans \qty{100}{\milli\litre} d'acide chlorhydrique à
\qty{1}{\mole\per\litre}.
\begin{enumerate}[label=\alph*)]
  \item Calcule les quantités de \ce{Zn} et de \ce{HCl}.
  \item Quel est le réactif limitant ?
  \item Quel volume de \ce{H2} (aux TPN) se dégage-t-il ?
\end{enumerate}
\end{exo}

\begin{exo}[avancement maximal et bilan]
Pour $\ce{2 A + 3 B -> C + 6 D}$, on part de $n_0(\ce{A})=\qty{0.10}{\mole}$ et
$n_0(\ce{B})=\qty{0.20}{\mole}$.
\begin{enumerate}[label=\alph*)]
  \item Calcule l'avancement maximal $\xi_{\max}=\min\!\big(n_i/\nu_i\big)$.
  \item Quel est le réactif limitant ?
  \item Donne le bilan final : quantités de \ce{A}, \ce{B}, \ce{C} et \ce{D}.
\end{enumerate}
\end{exo}

\begin{exo}[précipitation du sulfate de baryum]
On mélange \qty{200}{\milli\litre} de \ce{BaCl2} à \qty{0.5}{\mole\per\litre} avec
\qty{100}{\milli\litre} de \ce{Na2SO4} à \qty{0.5}{\mole\per\litre} :
$\ce{BaCl2 + Na2SO4 -> BaSO4 v + 2 NaCl}$. On donne $M(\ce{BaSO4})=\qty{233}{\gram\per\mole}$.
\begin{enumerate}[label=\alph*)]
  \item Calcule les quantités de \ce{BaCl2} et de \ce{Na2SO4}.
  \item Quel est le réactif limitant ?
  \item Quelle masse de \ce{BaSO4} se forme, et quelle quantité de \ce{BaCl2} reste-t-il ?
\end{enumerate}
\end{exo}

\end{document}
